\cotefile{156-6}
\scanpage{26}
\mspage{25}
D'abord une observation. Soit $F\in\unsure{\mathcal{F}|\mathscr{F}}$, d'où
\[
\mathrm{Omb}(F)=\lvert\unsure{\mathrm{Multomb}|\mathrm{Multiomb}}\,F\rvert
 \overset{\unsure{\mathrm{def}|\mathrm{déf}}}{=}
 \bigcup_{X\in\widetilde F}\deleted{\unsure{\widetilde X|X}}\mathrm{Omb}(\unsure{X|x})
\]
On a une application canonique
\[
\mathrm{Omb}(F)\xrightarrow{\varphi_F}\widetilde F
\]
\margin{NB. $\varphi_F$ induit $\mathrm{id}_{\widetilde F}$ sur $\widetilde F\subset\mathrm{Omb}(F)$}
associant à tt $Y\in\mathrm{Omb}\,F$, i.e. $Y\in M$ tel que
$\deleted{\unsure{Y\ll F|Y\in F}},\quad Y\ll F$, le plus petit \unsure{élément} de $\deleted{\unsure{G}}F^{\unsure{Y}}$
i.e. le plus petit $X$ de $\deleted{\unsure{G}}\widetilde F$ parmi ceux \unsure{satisfaisant}
$Y\ll X$ i.e. $\deleted{\unsure{O|\mathrm{Om}}}\,Y\in\mathrm{Omb}\,X$. Cette application est celle qui définit la relation de préordre de $\lvert\mathrm{Mulomb}(F)\rvert$ associée à la figure $\mathrm{Mulomb}(F)$.

Ceci dit, \inserted{pour $\deleted{\unsure{G}}F,\unsure{G|C}\in\mathcal{F}$ donnés,} dire que
$\underbrace{\deleted{\unsure{M}}\mathrm{Mulomb}\bigl(\deleted{\unsure{F}}G\bigr)}_{\Psi}\ll\underbrace{\unsure{\mathrm{Mulom}|\mathrm{Mulomb}}\,F}_{\Phi}$,
\unsure{signifie} que \deleted{\unsure{il|elle}}\unsure{l'on|on} a une inclusion
\[
 \lvert\Psi\rvert\subset\lvert\Phi\rvert\qquad\text{i.e.}\quad \mathrm{Om}\,G\subset\mathrm{Om}\,F,
\]
et une application \uline{croissante}
\[
\deleted{\unsure{\text{(1.42)}|\text{(1.47)}|\text{(1.48)}}}\qquad
f:\deleted{\unsure{\widetilde F}}\widetilde G\mathbin{\unsure{-|\longrightarrow}}\deleted{\unsure{\widetilde G}}\widetilde F
\]
rendant commutatif le diagramme
\begingroup\csname tagsleft@true\endcsname
\diagram{\[
\begin{tikzcd}[ampersand replacement=\&,column sep=large]
 \mathrm{Omb}\,\deleted{\unsure{F}}G \arrow[r,hook] \arrow[d,"\varphi_{\unsure{G|C}}"] \& \mathrm{Om}\,\deleted{\unsure{G}}F \arrow[d,"\varphi_F"]\\
 \deleted{\unsure{\widetilde F}}\widetilde G \arrow[r,"f" description] \& \widetilde F
\end{tikzcd}\quad .
\tag*{(1.48)}
\]}
\endgroup
\scanpage{27}
\mspage{26}
Montrons que cela équivaut à $G\ll F$. En effet, si $G\ll F$, associant à tout $X\in\widetilde G$ le plus petit élément $\deleted{\unsure{\varphi}}f(\unsure{x|X})$ de $\widetilde F$ tel que
\[
 X\ll f(x),
\]
on trouve bien une application
\[
 \deleted{\unsure{\widetilde G|G}}\qquad f:\widetilde G\longrightarrow\widetilde F
\]
et elle est \unsure{évidemment} croissante\,: \margin{si $X,X'\in\widetilde G$}
\[
\begin{aligned}
 X'\leq X\quad\text{i.e.}\quad X'\ll X\inserted{(\ll f(x))}\ &\text{implique}\quad X'\ll f(x)\\
 &\text{donc}\quad\deleted{\unsure{f}}\,f(x')\ll f(x)\\
 &\text{qui signifie}\quad f(x')\leq f(x)
\end{aligned}
\]
et il faut vérifier la commutativité de \unsure{(1,42)|(1.42)}, ce qui \deleted{\unsure{se voit comme à l'instant}.} \inserted{signifie}

\uline{Corollaire} Soient $X\in M$, $F,G\in\unsure{\mathcal F|\mathscr F}$ tels que $X\ll G\ll F$.
Soit $\deleted{\unsure{X|G}}X'$ le plus petit élément de $G^X$, $X''$ le plus petit élément de $F^{X'}$, alors $X''$ est \unsure{aussitt|aussi tt|aussi} le plus petit élément de $F^X$.

Il semble qu'il faille le poser en axiome, sous forme en apparence \unsure{affaiblie}\unsure{;|,} mais en fait équivalente. On va \deleted{\unsure{écrire|é}} \inserted{introduire} encore une notation, pour $X,Y\in M$
\begingroup\csname tagsleft@true\endcsname
\[
\begin{aligned}
 X\overset{\circ}{\ll}Y\qquad &(\unsure{-|=}X\text{ est \uline{\unsure{suffisamment} intérieur}}\\
 &\hspace{6em}\text{de }Y)
\end{aligned}
\tag*{(1.48)}
\]
\endgroup
pour
\[
 \deleted{X\ll Y,\ \text{et}\ X\ll Y}
\]
\[
 X\in\mathrm{Omb}(Y)^{\circ}=\mathrm{Omb}(Y)\setminus\bigcup_{Y'<Y}\mathrm{Omb}(Y')
\]
i.e. $X\ll Y$, et $X\ll Y'\leq Y$ implique $Y'=Y$.
\scanpage{28}
\mspage{27}
On doit donc poser

\uline{A\deleted{\unsure{C}}\unsure{E|t|C}7}\quad Si $X,Y,Z\in M$ sont tels que
\[
 X\overset{\circ}{\ll}Y\overset{\circ}{\ll}Z\qquad\text{alors}\qquad X\overset{\circ}{\ll}Z
\]
i.e. la relation $\overset{\circ}{\ll}$ est \uline{transitive}.

Ceci \unsure{assure} donc $\Rightarrow$ dans la prop.
En sens inverse, S'il \unsure{existe|exist} $f$ comme dans (1.42), soit $X\in\widetilde G$, \deleted{alors $\varphi_{\unsure{G|C}}(\unsure{X|x}$} \deleted{\unsure{on}} on sait que $\varphi_G(x)=X$ (tautologique\deleted{\unsure{?}}), moyennant (6), \deleted{on a alors} soit $f(x)=Y$, on a alors $Y=\varphi_F(x)$ i.e. $Y$ est le plus petit élément de $\widetilde F$ \deleted{\unsure{contenant} $X$}. \deleted{\unsure{Donc}} tel que $X\ll Y$. Cela \deleted{signifie} implique donc que pour tt $X\in\deleted{\unsure{\widetilde F}}\widetilde G$, \unsure{existe|exist} $Y\in\deleted{\unsure{\widetilde G}}\widetilde F$ telle que $X\ll Y$, i.e. $\deleted{\unsure{\widetilde G|\widetilde F}}$ (C4) $\widetilde G\ll\widetilde F$, \unsure{qed.|cqfd.} On n'a pas \unsure{utilisé|utilisé-} en fait que $f$ soit croissante, donc

\uline{Corollaire à la prop.}\quad \deleted{La condition,} \inserted{Pour} que $F\ll G$ i.e. $\unsure{\mathrm{Mulom}|\mathrm{Mulomb}}(F)\ll\mathrm{Mulomb}(G)$ il s. que $\mathrm{Omb}\,F\subset\mathrm{Omb}\,G$,
% slanted
\margin{Ce ne sera sans doute plus vrai, en \unsure{passant} aux \uline{petites} multifigures}
et que \deleted{\unsure{toute classe}} la relation d'équivalence \inserted{$\unsure{\mathcal R_F|R_F}$} sur $\mathrm{Omb}\,F$ définie par $\varphi_F$ soit plus fine que la relation d'équiv. induite \unsure{par|sur} $\mathrm{Omb}\,G$.
\scanpage{29}
\mspage{28}
\uline{Scholie}\quad Ainsi, via $F\longmapsto\mathrm{Mulomb}(F)$,
$(\unsure{\mathcal F|\mathscr F},\deleted{\unsure{\ll}}\leq,\ll)$ s'interprète comme un ``sous-atelier'' de $\mathrm{Fig}(M)$.

Donc la donnée de l'atelier $(\mathcal F,\leq,\ll)$ équivaut à celle d'un ens $M$, et d'une partie
\begingroup\csname tagsleft@true\endcsname
\[
 \deleted{\unsure{\Phi|\mathcal F}}\mathcal F_M\,
 \deleted{\unsure{\subset\mathcal P|\mathcal F}}\,
 \deleted{\Phi}\subset\mathrm{Fig}(M)\subset
 \unsure{\mathcal P|P}(\mathcal P(M))
 \tag*{(1.49)}
\]
\endgroup
\[
 \left(\text{donc}\quad
 \mathcal F_M\,\deleted{\unsure{\Phi}\,\unsure{\mathcal F_M|\Phi_M}}
 \in\mathcal P(\mathcal P(\mathcal P(M)))\right)
\]
satisfaisant certaines conditions \inserted{(qui doivent traduire C1 \unsure{;|à} C7)}.

Du coup, la notion de \inserted{(grande)} multiombre reprend un plus grand intérêt, et je me demande s'il ne faudrait pas une notation \deleted{\unsure{plus}} et un nom \unsure{pluss|plus} lapidaires.

Je voudrais expliciter les conditions sur cet ens de figures ensemblistes dans $M$. Je \unsure{désignerais|désignerai} par ${\unsure{\mathrm{At}|\mathrm{AE}}}^{*}1$ etc les axiomes en termes de $(M,\deleted{\unsure{\Phi_M}}\mathcal F_M\in\mathcal P\mathcal P\mathcal P(\unsure{m|M}))$.
\deleted{\unsure{$(\Phi_1)$|(1.49)|(1.50)}} Il est entendu que $\deleted{\unsure{\Phi}}\mathcal F_M\subset\mathrm{Fig}(M)$, i.e. \deleted{\uline{\unsure{At|AE}$^{*}$1}} que les éléments de $\deleted{\unsure{\Phi}}\mathcal F_M$ sont des ensembles de parties de $M$ qui sont des \uline{figures} ensemblistes $F$ dans $M$. \deleted{\illegible{1}} \deleted{Ainsi} Considérons
\scanpage{30}
\mspage{29}
\[
\begin{aligned}
 \overset{\deleted{\unsure{\sim|\approx}}}{M}^{*}
  &=\bigcup_{\deleted{\unsure{F}}\Phi\in\deleted{\unsure{\Phi}}\unsure{\mathcal F|\mathscr F}_M}
       \deleted{\unsure{F}}\Phi\subset\unsure{\mathcal P|P}(\deleted{\unsure{F}}M)\\
 &\quad\unsure{\shortparallel|\text{''}}\\
 &\left\{A\in\mathcal P(M)\mid\exists\,\deleted{\unsure{F}}\Phi\in\deleted{\unsure{\Phi}}\mathcal F_M
       \ \text{\unsure{avec}}\ A\in\deleted{\unsure{F}}\Phi\right\}
\end{aligned}
\]
i.e. l'ensemble des strates de toutes les figures $\in\deleted{\unsure{\Phi}.}\mathcal F_M$

\uline{At$^{*}$1}$_{\unsure{0|o}}$\quad $\overset{\deleted{\sim}}{M}^{*}$ est une figure ensembliste
\margin{\rule{0.45pt}{3.5\baselineskip}}
dans $M$, de support $M$. De plus, le préordre associé de $M$ \inserted{(soit $\ll$)} est une relation d'\uline{ordre} i.e.
\[
\begin{aligned}
 x,y\in M,\ x\neq y\quad\Longrightarrow\quad \exists\, A\in\overset{\deleted{\sim}}{M}^{*}
 &\quad\text{avec}\quad X\in A,\ Y\notin A\\
 &\quad\text{ou}\qquad Y\in A,\ x\notin A
\end{aligned}
\]

Ainsi, l'application
\[
 x\longmapsto \deleted{\unsure{M^{*}_{\ll x}|\widetilde M_{\ll x}}}M_{\ll x}
 \overset{\unsure{\mathrm{def}|\mathrm{déf}}}{=}\mathrm{Omb}(X)
\]
\begingroup\csname tagsleft@true\endcsname
\[
 M\longrightarrow\mathcal P(M)\tag*{(*)}
\]
\endgroup
induit une \uline{bijection} de $M$ sur $\overset{\deleted{\sim}}{M}^{*}$, transformant $\ll$ en la relation d'inclusion.
Les $A\in\overset{\deleted{\sim}}{M}^{*}$ sont les parties fermées \uline{irréductibles} de $M$ admettant un pt générique.

\deleted{\uline{\unsure{At|Al}$^{*}$2}}\quad Soit $\overset{\deleted{\sim}}{\deleted{\unsure{F}}\Phi}\in\deleted{\unsure{\Phi}}\inserted{\mathcal F_M}$, [\unsure{donc|\deleted{donc}} $\overset{\deleted{\sim}}{F}\subset\mathcal P(M)$ et par définition même $\deleted{\unsure{F}}\Phi\subset\overset{\deleted{\sim}}{M}^{*}\subset\mathcal P(M)$), donc $\Phi$ est
\scanpage{31}
\mspage{30}
l'image par \deleted{\unsure{(1.39)|(1.79)|(1.49)}} ($*$)
\[
 M\xrightarrow{\sim}\overset{\deleted{\unsure{\sim|\approx}}}{M}^{*}\subset\unsure{\mathcal P|P}(M)
\]
d'une unique \inserted{partie} $F\subset M$, de sorte que $\Phi=\widetilde F$, i.e.
\[
 \deleted{\unsure{\mathrm{At}|\mathrm{AC}}}\qquad
 \Phi=\bigl\{\underset{\inserted{\substack{\unsure{\shortparallel|\text{''}}\\ M_{\ll x}}}}
 {\mathrm{Omb}(\deleted{\unsure{\widetilde X}}X)}\mid X\in F\bigr\}
\]
Ainsi, on peut aussi interpréter, grâce à $\ll$, $\unsure{\mathcal F|\mathscr F}_M$ en termes d'un\deleted{e} \deleted{partie} ens de parties \inserted{$\widetilde{\mathcal F}$} de $M$
\diagram{\begin{tikzcd}[ampersand replacement=\&,column sep=large,row sep=1pt]
 \mathcal F_M \arrow[r,"\sim"] \& \widetilde{\mathcal F}\\
 \cap \& \cap\\
 \mathcal P(\mathcal P(M) \& \mathcal P(M)
\end{tikzcd}}

% slanted
\margin{cf GF $\overline{\underline{\unsure{\mathrm{IV}|\mathrm V}}}$ page 54 (Cor. 2) et page 55 (Remarque) Pour C3, cf page 58}
Ceci dit, \deleted{on demande que} \inserted{\unsure{Je veux|on veut|on veux} que $\widetilde{\mathcal F}\subset\mathcal P(M)$,}
\deleted{\unsure{soit associé à un}} \deleted{\unsure{satisfasse}} \inserted{\unsure{pour|muni|muni de}} \deleted{\unsure{une|cela|cette} relation d} la relation d'inclusion, satisfasse C1, C2, C3,
\deleted{et \unsure{identifiant|identifie} les \unsure{à}} \unsure{pour|que|par} $M$ \unsure{s'identifiant|s'identifie|s'identifier} à l'ens de ses objets \inserted{\unsure{irréd.|init.}}

\deleted{\uline{\unsure{At|AE}$^{*}$2} \unsure{bis}}\quad \inserted{On \unsure{voudra|veut donc} d'abord que} L'ens. \unsure{des|de} parties $\widetilde{\mathcal F}$ \deleted{de} $\unsure{\subset}\mathcal P(M)$ soit \uline{préadmissible}, i.e. pour tout $X\in M$, il existe un plus petit élément de $\widetilde{\mathcal F}$ qui le contienne. En d'autres termes,

\uline{At$^{*}$2}$_{\unsure{0|o}}$\quad
\margin{\rule{0.45pt}{4\baselineskip}}
\inserted{Pour} \unsure{tout|toute} $\deleted{\Phi\in\mathcal F_M}$ $M$-figure ensembliste $\Phi$ \deleted{\unsure{de|dans}} $\in\mathcal F_M$, et pour $A$ [$=\mathrm{Omb}(x)$!] \deleted{\unsure{de}} \unsure{strate} de $\Phi$ (i.e. $A\in\Phi$), il existe une \uline{plus petite figure} \inserted{$\Phi_A$} \deleted{\unsure{dans|tous} \unsure{des}} $\in\overset{\deleted{\sim}}{\mathcal F}_M$
% slanted
\margin{c'est la \uline{\unsure{multistrate}} définie par $A$}
(plus petite au sens de l'\uline{inclusion} \deleted{\unsure{)}} entre figures i.e. entre parties de $\mathcal P(M)$, ou \deleted{de} $M^{*}\subset\mathcal P(M)$) qui contienne $A$ comme élément (laquelle $\Phi_A$ sera donc $\subset\Phi$).

% slanted
\margin{C'est C1 pour $\unsure{\mathcal F_M|F_M}$}
\uline{At$^{*}$3}$_0$\quad Si $\Phi\in\mathcal F_M$, et si \deleted{\unsure{des}} $\Phi_i\in\mathcal F_M$, $\Phi_i\subset\Phi$, alors $\bigcup\Phi_i\in\mathcal F_M$ (stabilité par réunions \uline{conjointes} \inserted{dans $\mathcal F_M$})
\scanpage{32}
\mspage{31}
et de plus on a $\varnothing\in\unsure{\mathcal F|\mathscr F}_M$ (ou encore $\mathcal F_M\neq\varnothing$).

Enfin, on voudra aussi

\uline{At$^{*}$4}$_{\unsure{0|o}}$\quad Si $\Phi,\Psi\in\deleted{\unsure{\Phi}}\mathcal F_M$, pour que $\Phi\cup\Psi$ soit \deleted{élément} dans $\mathcal F_M$ \inserted{(i.e., par At$^{*}$3, $\Phi,\Psi$ \unsure{majorés} dans $\mathcal F_M$ \inserted{pour l'inclusion})}, il f. et s. que pour tout $A\in\Phi$, $\unsure{\cdot|.}B\in\Psi$, $\Phi_A\cup\Phi_B$ soit dans $\mathcal F_M$.

Je voudrais revenir sur ce $\Phi_A$. En fait, $\Phi_A$ consiste en l'ensemble des sous-strates de $A$ dans $\Phi$, donc \uline{At$^{*}$2}$_0$ peut se formuler sous la forme plus compréhensible

\uline{At$^{*}$}2$_0$\quad Soient $\Phi$ et $\Psi$ \deleted{\unsure{dans}}\unsure{deux} $M$-figures $\in\mathcal F_M$, et soit $A\in\Phi\cap\Psi$ une strate commune. Alors la multistrate associée à $A$ dans $\Phi$, et dans $\Psi$, est la \deleted{\unsure{m}}\unsure{même}, et elle est $\in\unsure{,|.}\mathcal F_M$.

\uline{Corollaire}\quad Si $\Phi,\Psi\in\mathcal F_M$ avec $\Phi\subset\Psi$, alors $\Phi$ est \unsure{fermée|fermé} dans $\Psi$, donc est une sous figure $(\Phi\leq\Psi)$.
\scanpage{33}
\mspage{32}
[Il s'ensuit qu'on a \unsure{dans|donc} $M\simeq M^{*}$ une relation d'ordre $\leq$, avec
\[
\begin{aligned}
 X\leq Y\quad\Longleftrightarrow\quad&
 \deleted{\mathrm{Omb}(x)}\Phi_{\mathrm{Omb}(x)}
 \subset\Phi_{\mathrm{Omb}(Y)}\\
 &\qquad\deleted{\text{i.e.}}\ \text{i.e.}\quad
 \mathrm{Omb}(x)\in
 \underset{\inserted{\substack{\unsure{\mathrm{def}|\mathrm{déf}}\ \unsure{\shortparallel|\text{''}}\\
 \unsure{\mathrm{Mulomb}|\mathrm{Multomb}}(Y)}}}{\Phi_{\mathrm{Omb}(Y)}},
\end{aligned}
\]
de telle façon que les \deleted{éléments de} $F\in\widetilde{\unsure{\mathcal F|\mathscr F}}$ soient des parties fermées de $M$.
Associant à tout $X\in M$ la partie $\deleted{\unsure{\widetilde X|X}}$
\[
\begin{aligned}
 \widetilde X&=\{Y\in M\mid Y\leq\unsure{x|X}\ \text{i.e.}\quad
 \mathrm{Omb}(Y)\in\Phi_{\mathrm{Omb}(x)}\}\\
 &=M_{\leq x}
\end{aligned}
\]
on trouve une application injective
\[
 M\hookrightarrow\widetilde{\mathcal F}\quad(\simeq\mathcal F_M)
\]
qui induit un iso entre $M$ et l'ens des éléments str. irréductibles de $\widetilde{\mathcal F}$ (ordonné par inclusion, notée $\leq$).]

\deleted{\unsure{Donc}}

\uline{Théorème} - \unsure{\uline{Scholie}|\uuline{Scholie}}\quad La catégorie \inserted{(groupoïde)} des \deleted{\unsure{$\mathcal F$}} ateliers $(\mathcal F,\leq,\ll)$ satisfaisant à \deleted{\unsure{C}}At 1 \unsure{.|à|,} At 7, équivaut à celle des $(M,\mathcal F_M)$ (avec $\mathcal F_M\in\deleted{\mathcal P\mathcal P}\mathrm{Fig}(M)$) satisfaisant les conditions suivantes

\uline{At$^{*}$1}, $\mathcal F_M$ \deleted{non vide, et} \inserted{contient la figure vide, \unsure{et|est}} stable par \deleted{\unsure{les sup majorés|sup majorés}} \inserted{réunions} majorées (dans $\mathcal F_M$)

\deleted{\unsure{Cela|Ceci}} (La stabilité par sup \unsure{équivaudra|équivaudre}, moyennant la condition suivante, à celle de \deleted{\unsure{passage}} stabilité par passage à une \uline{sous} \deleted{\unsure{famille fermée|famille finie}} \uline{figure}.)
